% =============================================================================
%  impact.tex -- Section: Projected Impact Assessment
%  BodyCheck -- AI-Powered Medical Triage & 3D Symptom Checker
%  Vietnamese Student HackAIthon 2026 -- Bang B (Challenger)
%  Team Clavicular | HCMUS APCS/CLC
%
%  Pure content module -- include via % =============================================================================
%  impact.tex -- Section: Projected Impact Assessment
%  BodyCheck -- AI-Powered Medical Triage & 3D Symptom Checker
%  Vietnamese Student HackAIthon 2026 -- Bang B (Challenger)
%  Team Clavicular | HCMUS APCS/CLC
%
%  Pure content module -- include via % =============================================================================
%  impact.tex -- Section: Projected Impact Assessment
%  BodyCheck -- AI-Powered Medical Triage & 3D Symptom Checker
%  Vietnamese Student HackAIthon 2026 -- Bang B (Challenger)
%  Team Clavicular | HCMUS APCS/CLC
%
%  Pure content module -- include via % =============================================================================
%  impact.tex -- Section: Projected Impact Assessment
%  BodyCheck -- AI-Powered Medical Triage & 3D Symptom Checker
%  Vietnamese Student HackAIthon 2026 -- Bang B (Challenger)
%  Team Clavicular | HCMUS APCS/CLC
%
%  Pure content module -- include via \input{impact}.
%  Uses only packages and colour definitions already loaded in main.tex.
% =============================================================================

\section{Projected Impact Assessment}
\label{sec:impact}

This section addresses the fourth Round~1 evaluation criterion ---
Projected Impact. We evaluate BodyCheck along four dimensions required by
the rubric: social and business value; potential user scale using a
TAM--SAM--SOM framing; competitive advantage in the target market; and the
revenue/value model that can sustain the product after the competition.

The impact thesis is straightforward: BodyCheck does not try to replace
doctors. Its highest-value contribution is to improve the earliest minutes of
a health concern, when users are uncertain, anxious, and often forced to choose
between ignoring symptoms, searching online, or visiting a clinic without a
clear sense of urgency. By turning vague symptoms into structured triage
guidance, BodyCheck can reduce confusion for users, prepare cleaner intake
summaries for healthcare teams, and create a reusable learning environment for
health-literacy and medical-training use cases.

Table~\ref{tab:impact_matrix} summarises the expected impact channels.

% ---------------------------------------------------------------------------
% Table 1 -- Projected Impact Matrix
% ---------------------------------------------------------------------------
\begin{table}[htbp]
  \centering
  \caption{Projected impact channels for BodyCheck aligned to the
    HackAIthon 2026 Round~1 impact rubric.}
  \label{tab:impact_matrix}
  \setlength{\tabcolsep}{5pt}
  \renewcommand{\arraystretch}{1.45}
  \small
  \begin{tabular}{@{} p{3.0cm} p{4.6cm} p{4.7cm} @{}}
    \toprule
    \textbf{Impact Area} &
    \textbf{Expected Benefit} &
    \textbf{Measurement Signal} \\
    \midrule

    Social Benefit &
    Earlier recognition of red flags, clearer self-care guidance, and
    faster escalation to professional care when symptoms appear risky. &
    Time-to-first-guidance, emergency-escalation rate, user clarity score,
    and clinic handoff clicks. \\

    Healthcare Operations &
    Cleaner pre-consultation summaries reduce repetitive intake work for
    triage operators, telehealth nurses, and primary-care clinics. &
    Intake completion rate, summary reuse rate, average number of follow-up
    clarification questions. \\

    Education &
    Interactive 3D body-region selection and AI-generated triage cases create
    a simulation tool for medical, nursing, first-aid, and health-literacy
    programmes. &
    Number of institutional pilots, completed learning cases, instructor
    feedback score. \\

    Business Value &
    One shared technical core supports three commercial channels: consumer
    app, healthcare API, and education licence. &
    Monthly active users, API request volume, institutional pilots, paid
    conversion rate. \\

    \bottomrule
  \end{tabular}
\end{table}

% =============================================================================
\subsection{Social and Healthcare Impact}
\label{sec:impact-social}
% =============================================================================

BodyCheck creates value by reducing the gap between symptom onset and an
appropriate next action. For non-specialist users, the product translates a
stressful and ambiguous experience --- for example, ``chest tightness'',
``shoulder pain'', or ``lower abdominal discomfort'' --- into a structured
triage response:

\begin{itemize}[leftmargin=2em, itemsep=3pt, topsep=4pt]
  \item possible explanations expressed with uncertainty language;
  \item risk level and recommended timing for care;
  \item red-flag symptoms that require urgent escalation;
  \item home-care tips for low-risk situations;
  \item trusted references and nearby care options when professional help is
        appropriate.
\end{itemize}

This matters because the first decision is often the most confusing one:
whether to wait, self-care, schedule a routine appointment, or seek urgent
care. BodyCheck improves that decision without making a final diagnosis. The
AI response is deliberately constrained by prompt rules, a structured schema,
and a medical disclaimer so that the user receives triage support rather than
false diagnostic certainty.

For healthcare operations, the impact is not only patient-facing. A structured
intake flow can produce cleaner summaries for telehealth nurses, clinic
reception teams, and primary-care doctors. Instead of receiving a long
unstructured complaint, the clinician can see the affected region, duration,
severity, possible red flags, AI-generated follow-up questions, and the user's
own symptom notes. This can reduce repeated clarification work and make the
first human consultation more focused.

For education, the same workflow becomes a learning tool. Students can explore
anatomical regions on a 3D model, compare symptoms across body systems, and
observe how different risk factors change the recommended triage pathway. This
creates a second impact channel beyond consumer health: interactive clinical
reasoning practice for medical and nursing students, first-aid training
programmes, and public health-literacy courses.

% =============================================================================
\subsection{Market Potential and TAM--SAM--SOM Framing}
\label{sec:impact-market}
% =============================================================================

At Round~1 stage, exact commercial market sizing is less reliable than a
transparent bottom-up model. BodyCheck therefore uses a conservative
TAM--SAM--SOM framing based on reachable digital users, expected symptom-help
seeking behaviour, and the team's realistic first-year distribution channels.

\begin{table}[htbp]
  \centering
  \caption{Conservative TAM--SAM--SOM framing for BodyCheck. Values are
    planning estimates for product direction, not claimed realised traction.}
  \label{tab:tam_sam_som}
  \setlength{\tabcolsep}{5pt}
  \renewcommand{\arraystretch}{1.45}
  \small
  \begin{tabular}{@{} p{2.0cm} p{5.2cm} p{5.1cm} @{}}
    \toprule
    \textbf{Layer} & \textbf{Definition} & \textbf{BodyCheck Estimate} \\
    \midrule

    TAM &
    Digitally reachable users who may need symptom guidance, health
    information, or basic triage support. &
    Domestic floor: Vietnam's internet-user base already identified in
    \S\ref{sec:problem-fit}. Global upside comes from the English-first
    interface and web/PWA distribution. \\

    SAM &
    Users and organisations reachable by the initial product scope: English
    and Vietnamese web/mobile users, telehealth intake teams, clinics, and
    education providers. &
    Initial serviceable segment: students, young families, health-conscious
    consumers, HCMUS/HCMC beta networks, online first-aid communities, and
    small healthcare/education partners. \\

    SOM &
    Share realistically obtainable in the first 12 months after the
    competition, assuming a student-led team and low marketing budget. &
    Target: 30--50 closed-beta testers, 10{,}000--30{,}000 first-year
    consumer users, 2--3 education pilots, and 1--2 healthcare workflow
    pilots by June~2027. \\

    \bottomrule
  \end{tabular}
\end{table}

The strongest initial market is not the full healthcare system. It is the
pre-clinical decision layer: people who feel symptoms and need immediate,
understandable next-step guidance before they reach a professional. This layer
is large because it sits upstream of clinics, hospitals, telehealth platforms,
and search engines. BodyCheck's web-first architecture lets the team test this
market with a low-friction PWA before investing in heavier clinical
integrations.

The education segment provides an additional serviceable market with lower
regulatory friction. A 3D anatomical symptom checker can be packaged as a
simulation environment for instructors, allowing students to practise
triage-style reasoning without involving real patient data. This makes the
product useful even before it becomes a clinically reviewed public health tool.

% =============================================================================
\subsection{Competitive Advantage}
\label{sec:impact-competitive}
% =============================================================================

BodyCheck's competitive advantage comes from combining several capabilities
that usually appear separately in existing tools. Search engines provide
information but no structured triage path. Static symptom checkers provide
forms but often feel rigid and text-heavy. Generic chatbots can converse, but
they do not know where the user is pointing on the body and are not designed
around care-navigation handoff. Telehealth platforms connect users to
clinicians, but they often begin after the user has already decided to seek
care.

BodyCheck differs by making the body region the first-class input. The user can
click or paint affected areas on an interactive 3D anatomical model before
writing symptoms. This reduces ambiguity for non-specialist users and creates a
more structured clinical context for the AI layer. The resulting response is
not just a chat answer; it is a triage packet containing body region, symptom
summary, risk level, possible conditions, warning signs, references, and nearby
care options.

\begin{table}[htbp]
  \centering
  \caption{BodyCheck differentiation against common alternatives.}
  \label{tab:competitive_comparison}
  \setlength{\tabcolsep}{5pt}
  \renewcommand{\arraystretch}{1.45}
  \small
  \begin{tabular}{@{} p{3.2cm} p{4.8cm} p{4.8cm} @{}}
    \toprule
    \textbf{Alternative} & \textbf{Limitation} & \textbf{BodyCheck Difference} \\
    \midrule

    Search Engines &
    Return many documents; users must interpret urgency themselves. &
    Produces a constrained triage summary with next-step timing and red-flag
    escalation. \\

    Static Symptom Checkers &
    Depend on rigid forms and predefined symptom trees. &
    Supports natural language, voice input, and targeted follow-up questions. \\

    Generic AI Chatbots &
    Lack anatomical context, care-routing workflow, and product-level safety
    framing. &
    Combines 3D body-region input, structured diagnosis schema, references,
    disclaimer, and clinic handoff. \\

    Telehealth Intake Forms &
    Often start after the user has already chosen to seek care. &
    Supports the earlier decision layer and can still export structured
    summaries into telehealth workflows. \\

    \bottomrule
  \end{tabular}
\end{table}

This combination creates more than a cosmetic difference. The core workflow is
at least 30\% different from conventional text-only symptom tools because the
product changes the input modality, response structure, and downstream action
path. It begins with spatial symptom localisation, adds LLM-based reasoning,
preserves safety constraints, and ends with either self-care guidance,
professional escalation, or a clinic/education handoff.

% =============================================================================
\subsection{Revenue and Value Model}
\label{sec:impact-revenue}
% =============================================================================

BodyCheck's revenue model is designed around three complementary channels so
that the product is not dependent on a single customer type.

\begin{enumerate}[leftmargin=2em, itemsep=4pt, topsep=4pt]
  \item \textbf{B2C freemium consumer app:}
        Basic triage guidance remains free to maximise social reach. Paid
        features can include saved symptom history, voice accessibility,
        advanced source summaries, care-navigation reminders, and family
        profile management.
  \item \textbf{B2B2C healthcare API licensing:}
        Telehealth operators, clinics, hospital intake portals, and insurance
        triage teams can license BodyCheck as a white-label REST API. Pricing
        combines a platform fee with per-request usage, following the same
        linear cost logic defined in the feasibility section.
  \item \textbf{Education institutional licence:}
        Universities, medical training programmes, first-aid organisations,
        and online learning platforms can use BodyCheck as a white-label
        simulation tool. This channel is priced as an annual licence because
        value is driven by curriculum usage rather than individual medical
        sessions.
\end{enumerate}

The value model differs by channel. For consumers, the value is clarity and
confidence in the next step. For healthcare partners, the value is a more
structured pre-consultation intake flow and lower repetitive clarification
work. For education partners, the value is interactive, repeatable anatomy and
triage training without exposing real patient data.

% =============================================================================
\subsection{Impact Measurement Plan}
\label{sec:impact-measurement}
% =============================================================================

The post-competition roadmap in \S\ref{sec:feasibility-gtm} becomes stronger
when paired with measurable impact indicators. BodyCheck will track impact
through anonymised product and workflow metrics rather than raw symptom logs:

\begin{itemize}[leftmargin=2em, itemsep=3pt, topsep=4pt]
  \item \textbf{Access:} number of completed triage sessions, active users,
        returning users, and device/platform distribution.
  \item \textbf{Clarity:} post-session rating on whether the user understood
        what to do next.
  \item \textbf{Safety:} number of red-flag escalations, emergency-care
        recommendation displays, and disclaimer visibility checks.
  \item \textbf{Workflow value:} number of generated intake summaries,
        clinician/partner feedback on summary usefulness, and reduction in
        repeated clarification questions during pilot workflows.
  \item \textbf{Education value:} number of completed training cases,
        instructor feedback, and learner quiz-score improvement before and
        after using the simulation flow.
\end{itemize}

These metrics allow the team to demonstrate social benefit without making
unverifiable medical-outcome claims. The goal for the first 12 months is not
to prove that BodyCheck replaces clinical care. The goal is to prove that the
product helps users understand urgency earlier, helps care teams receive
cleaner symptom context, and helps educators teach anatomical symptom
recognition in a more interactive way.

% =============================================================================
% End of Section: Projected Impact Assessment
% =============================================================================
.
%  Uses only packages and colour definitions already loaded in main.tex.
% =============================================================================

\section{Projected Impact Assessment}
\label{sec:impact}

This section addresses the fourth Round~1 evaluation criterion ---
Projected Impact. We evaluate BodyCheck along four dimensions required by
the rubric: social and business value; potential user scale using a
TAM--SAM--SOM framing; competitive advantage in the target market; and the
revenue/value model that can sustain the product after the competition.

The impact thesis is straightforward: BodyCheck does not try to replace
doctors. Its highest-value contribution is to improve the earliest minutes of
a health concern, when users are uncertain, anxious, and often forced to choose
between ignoring symptoms, searching online, or visiting a clinic without a
clear sense of urgency. By turning vague symptoms into structured triage
guidance, BodyCheck can reduce confusion for users, prepare cleaner intake
summaries for healthcare teams, and create a reusable learning environment for
health-literacy and medical-training use cases.

Table~\ref{tab:impact_matrix} summarises the expected impact channels.

% ---------------------------------------------------------------------------
% Table 1 -- Projected Impact Matrix
% ---------------------------------------------------------------------------
\begin{table}[htbp]
  \centering
  \caption{Projected impact channels for BodyCheck aligned to the
    HackAIthon 2026 Round~1 impact rubric.}
  \label{tab:impact_matrix}
  \setlength{\tabcolsep}{5pt}
  \renewcommand{\arraystretch}{1.45}
  \small
  \begin{tabular}{@{} p{3.0cm} p{4.6cm} p{4.7cm} @{}}
    \toprule
    \textbf{Impact Area} &
    \textbf{Expected Benefit} &
    \textbf{Measurement Signal} \\
    \midrule

    Social Benefit &
    Earlier recognition of red flags, clearer self-care guidance, and
    faster escalation to professional care when symptoms appear risky. &
    Time-to-first-guidance, emergency-escalation rate, user clarity score,
    and clinic handoff clicks. \\

    Healthcare Operations &
    Cleaner pre-consultation summaries reduce repetitive intake work for
    triage operators, telehealth nurses, and primary-care clinics. &
    Intake completion rate, summary reuse rate, average number of follow-up
    clarification questions. \\

    Education &
    Interactive 3D body-region selection and AI-generated triage cases create
    a simulation tool for medical, nursing, first-aid, and health-literacy
    programmes. &
    Number of institutional pilots, completed learning cases, instructor
    feedback score. \\

    Business Value &
    One shared technical core supports three commercial channels: consumer
    app, healthcare API, and education licence. &
    Monthly active users, API request volume, institutional pilots, paid
    conversion rate. \\

    \bottomrule
  \end{tabular}
\end{table}

% =============================================================================
\subsection{Social and Healthcare Impact}
\label{sec:impact-social}
% =============================================================================

BodyCheck creates value by reducing the gap between symptom onset and an
appropriate next action. For non-specialist users, the product translates a
stressful and ambiguous experience --- for example, ``chest tightness'',
``shoulder pain'', or ``lower abdominal discomfort'' --- into a structured
triage response:

\begin{itemize}[leftmargin=2em, itemsep=3pt, topsep=4pt]
  \item possible explanations expressed with uncertainty language;
  \item risk level and recommended timing for care;
  \item red-flag symptoms that require urgent escalation;
  \item home-care tips for low-risk situations;
  \item trusted references and nearby care options when professional help is
        appropriate.
\end{itemize}

This matters because the first decision is often the most confusing one:
whether to wait, self-care, schedule a routine appointment, or seek urgent
care. BodyCheck improves that decision without making a final diagnosis. The
AI response is deliberately constrained by prompt rules, a structured schema,
and a medical disclaimer so that the user receives triage support rather than
false diagnostic certainty.

For healthcare operations, the impact is not only patient-facing. A structured
intake flow can produce cleaner summaries for telehealth nurses, clinic
reception teams, and primary-care doctors. Instead of receiving a long
unstructured complaint, the clinician can see the affected region, duration,
severity, possible red flags, AI-generated follow-up questions, and the user's
own symptom notes. This can reduce repeated clarification work and make the
first human consultation more focused.

For education, the same workflow becomes a learning tool. Students can explore
anatomical regions on a 3D model, compare symptoms across body systems, and
observe how different risk factors change the recommended triage pathway. This
creates a second impact channel beyond consumer health: interactive clinical
reasoning practice for medical and nursing students, first-aid training
programmes, and public health-literacy courses.

% =============================================================================
\subsection{Market Potential and TAM--SAM--SOM Framing}
\label{sec:impact-market}
% =============================================================================

At Round~1 stage, exact commercial market sizing is less reliable than a
transparent bottom-up model. BodyCheck therefore uses a conservative
TAM--SAM--SOM framing based on reachable digital users, expected symptom-help
seeking behaviour, and the team's realistic first-year distribution channels.

\begin{table}[htbp]
  \centering
  \caption{Conservative TAM--SAM--SOM framing for BodyCheck. Values are
    planning estimates for product direction, not claimed realised traction.}
  \label{tab:tam_sam_som}
  \setlength{\tabcolsep}{5pt}
  \renewcommand{\arraystretch}{1.45}
  \small
  \begin{tabular}{@{} p{2.0cm} p{5.2cm} p{5.1cm} @{}}
    \toprule
    \textbf{Layer} & \textbf{Definition} & \textbf{BodyCheck Estimate} \\
    \midrule

    TAM &
    Digitally reachable users who may need symptom guidance, health
    information, or basic triage support. &
    Domestic floor: Vietnam's internet-user base already identified in
    \S\ref{sec:problem-fit}. Global upside comes from the English-first
    interface and web/PWA distribution. \\

    SAM &
    Users and organisations reachable by the initial product scope: English
    and Vietnamese web/mobile users, telehealth intake teams, clinics, and
    education providers. &
    Initial serviceable segment: students, young families, health-conscious
    consumers, HCMUS/HCMC beta networks, online first-aid communities, and
    small healthcare/education partners. \\

    SOM &
    Share realistically obtainable in the first 12 months after the
    competition, assuming a student-led team and low marketing budget. &
    Target: 30--50 closed-beta testers, 10{,}000--30{,}000 first-year
    consumer users, 2--3 education pilots, and 1--2 healthcare workflow
    pilots by June~2027. \\

    \bottomrule
  \end{tabular}
\end{table}

The strongest initial market is not the full healthcare system. It is the
pre-clinical decision layer: people who feel symptoms and need immediate,
understandable next-step guidance before they reach a professional. This layer
is large because it sits upstream of clinics, hospitals, telehealth platforms,
and search engines. BodyCheck's web-first architecture lets the team test this
market with a low-friction PWA before investing in heavier clinical
integrations.

The education segment provides an additional serviceable market with lower
regulatory friction. A 3D anatomical symptom checker can be packaged as a
simulation environment for instructors, allowing students to practise
triage-style reasoning without involving real patient data. This makes the
product useful even before it becomes a clinically reviewed public health tool.

% =============================================================================
\subsection{Competitive Advantage}
\label{sec:impact-competitive}
% =============================================================================

BodyCheck's competitive advantage comes from combining several capabilities
that usually appear separately in existing tools. Search engines provide
information but no structured triage path. Static symptom checkers provide
forms but often feel rigid and text-heavy. Generic chatbots can converse, but
they do not know where the user is pointing on the body and are not designed
around care-navigation handoff. Telehealth platforms connect users to
clinicians, but they often begin after the user has already decided to seek
care.

BodyCheck differs by making the body region the first-class input. The user can
click or paint affected areas on an interactive 3D anatomical model before
writing symptoms. This reduces ambiguity for non-specialist users and creates a
more structured clinical context for the AI layer. The resulting response is
not just a chat answer; it is a triage packet containing body region, symptom
summary, risk level, possible conditions, warning signs, references, and nearby
care options.

\begin{table}[htbp]
  \centering
  \caption{BodyCheck differentiation against common alternatives.}
  \label{tab:competitive_comparison}
  \setlength{\tabcolsep}{5pt}
  \renewcommand{\arraystretch}{1.45}
  \small
  \begin{tabular}{@{} p{3.2cm} p{4.8cm} p{4.8cm} @{}}
    \toprule
    \textbf{Alternative} & \textbf{Limitation} & \textbf{BodyCheck Difference} \\
    \midrule

    Search Engines &
    Return many documents; users must interpret urgency themselves. &
    Produces a constrained triage summary with next-step timing and red-flag
    escalation. \\

    Static Symptom Checkers &
    Depend on rigid forms and predefined symptom trees. &
    Supports natural language, voice input, and targeted follow-up questions. \\

    Generic AI Chatbots &
    Lack anatomical context, care-routing workflow, and product-level safety
    framing. &
    Combines 3D body-region input, structured diagnosis schema, references,
    disclaimer, and clinic handoff. \\

    Telehealth Intake Forms &
    Often start after the user has already chosen to seek care. &
    Supports the earlier decision layer and can still export structured
    summaries into telehealth workflows. \\

    \bottomrule
  \end{tabular}
\end{table}

This combination creates more than a cosmetic difference. The core workflow is
at least 30\% different from conventional text-only symptom tools because the
product changes the input modality, response structure, and downstream action
path. It begins with spatial symptom localisation, adds LLM-based reasoning,
preserves safety constraints, and ends with either self-care guidance,
professional escalation, or a clinic/education handoff.

% =============================================================================
\subsection{Revenue and Value Model}
\label{sec:impact-revenue}
% =============================================================================

BodyCheck's revenue model is designed around three complementary channels so
that the product is not dependent on a single customer type.

\begin{enumerate}[leftmargin=2em, itemsep=4pt, topsep=4pt]
  \item \textbf{B2C freemium consumer app:}
        Basic triage guidance remains free to maximise social reach. Paid
        features can include saved symptom history, voice accessibility,
        advanced source summaries, care-navigation reminders, and family
        profile management.
  \item \textbf{B2B2C healthcare API licensing:}
        Telehealth operators, clinics, hospital intake portals, and insurance
        triage teams can license BodyCheck as a white-label REST API. Pricing
        combines a platform fee with per-request usage, following the same
        linear cost logic defined in the feasibility section.
  \item \textbf{Education institutional licence:}
        Universities, medical training programmes, first-aid organisations,
        and online learning platforms can use BodyCheck as a white-label
        simulation tool. This channel is priced as an annual licence because
        value is driven by curriculum usage rather than individual medical
        sessions.
\end{enumerate}

The value model differs by channel. For consumers, the value is clarity and
confidence in the next step. For healthcare partners, the value is a more
structured pre-consultation intake flow and lower repetitive clarification
work. For education partners, the value is interactive, repeatable anatomy and
triage training without exposing real patient data.

% =============================================================================
\subsection{Impact Measurement Plan}
\label{sec:impact-measurement}
% =============================================================================

The post-competition roadmap in \S\ref{sec:feasibility-gtm} becomes stronger
when paired with measurable impact indicators. BodyCheck will track impact
through anonymised product and workflow metrics rather than raw symptom logs:

\begin{itemize}[leftmargin=2em, itemsep=3pt, topsep=4pt]
  \item \textbf{Access:} number of completed triage sessions, active users,
        returning users, and device/platform distribution.
  \item \textbf{Clarity:} post-session rating on whether the user understood
        what to do next.
  \item \textbf{Safety:} number of red-flag escalations, emergency-care
        recommendation displays, and disclaimer visibility checks.
  \item \textbf{Workflow value:} number of generated intake summaries,
        clinician/partner feedback on summary usefulness, and reduction in
        repeated clarification questions during pilot workflows.
  \item \textbf{Education value:} number of completed training cases,
        instructor feedback, and learner quiz-score improvement before and
        after using the simulation flow.
\end{itemize}

These metrics allow the team to demonstrate social benefit without making
unverifiable medical-outcome claims. The goal for the first 12 months is not
to prove that BodyCheck replaces clinical care. The goal is to prove that the
product helps users understand urgency earlier, helps care teams receive
cleaner symptom context, and helps educators teach anatomical symptom
recognition in a more interactive way.

% =============================================================================
% End of Section: Projected Impact Assessment
% =============================================================================
.
%  Uses only packages and colour definitions already loaded in main.tex.
% =============================================================================

\section{Projected Impact Assessment}
\label{sec:impact}

This section addresses the fourth Round~1 evaluation criterion ---
Projected Impact. We evaluate BodyCheck along four dimensions required by
the rubric: social and business value; potential user scale using a
TAM--SAM--SOM framing; competitive advantage in the target market; and the
revenue/value model that can sustain the product after the competition.

The impact thesis is straightforward: BodyCheck does not try to replace
doctors. Its highest-value contribution is to improve the earliest minutes of
a health concern, when users are uncertain, anxious, and often forced to choose
between ignoring symptoms, searching online, or visiting a clinic without a
clear sense of urgency. By turning vague symptoms into structured triage
guidance, BodyCheck can reduce confusion for users, prepare cleaner intake
summaries for healthcare teams, and create a reusable learning environment for
health-literacy and medical-training use cases.

Table~\ref{tab:impact_matrix} summarises the expected impact channels.

% ---------------------------------------------------------------------------
% Table 1 -- Projected Impact Matrix
% ---------------------------------------------------------------------------
\begin{table}[htbp]
  \centering
  \caption{Projected impact channels for BodyCheck aligned to the
    HackAIthon 2026 Round~1 impact rubric.}
  \label{tab:impact_matrix}
  \setlength{\tabcolsep}{5pt}
  \renewcommand{\arraystretch}{1.45}
  \small
  \begin{tabular}{@{} p{3.0cm} p{4.6cm} p{4.7cm} @{}}
    \toprule
    \textbf{Impact Area} &
    \textbf{Expected Benefit} &
    \textbf{Measurement Signal} \\
    \midrule

    Social Benefit &
    Earlier recognition of red flags, clearer self-care guidance, and
    faster escalation to professional care when symptoms appear risky. &
    Time-to-first-guidance, emergency-escalation rate, user clarity score,
    and clinic handoff clicks. \\

    Healthcare Operations &
    Cleaner pre-consultation summaries reduce repetitive intake work for
    triage operators, telehealth nurses, and primary-care clinics. &
    Intake completion rate, summary reuse rate, average number of follow-up
    clarification questions. \\

    Education &
    Interactive 3D body-region selection and AI-generated triage cases create
    a simulation tool for medical, nursing, first-aid, and health-literacy
    programmes. &
    Number of institutional pilots, completed learning cases, instructor
    feedback score. \\

    Business Value &
    One shared technical core supports three commercial channels: consumer
    app, healthcare API, and education licence. &
    Monthly active users, API request volume, institutional pilots, paid
    conversion rate. \\

    \bottomrule
  \end{tabular}
\end{table}

% =============================================================================
\subsection{Social and Healthcare Impact}
\label{sec:impact-social}
% =============================================================================

BodyCheck creates value by reducing the gap between symptom onset and an
appropriate next action. For non-specialist users, the product translates a
stressful and ambiguous experience --- for example, ``chest tightness'',
``shoulder pain'', or ``lower abdominal discomfort'' --- into a structured
triage response:

\begin{itemize}[leftmargin=2em, itemsep=3pt, topsep=4pt]
  \item possible explanations expressed with uncertainty language;
  \item risk level and recommended timing for care;
  \item red-flag symptoms that require urgent escalation;
  \item home-care tips for low-risk situations;
  \item trusted references and nearby care options when professional help is
        appropriate.
\end{itemize}

This matters because the first decision is often the most confusing one:
whether to wait, self-care, schedule a routine appointment, or seek urgent
care. BodyCheck improves that decision without making a final diagnosis. The
AI response is deliberately constrained by prompt rules, a structured schema,
and a medical disclaimer so that the user receives triage support rather than
false diagnostic certainty.

For healthcare operations, the impact is not only patient-facing. A structured
intake flow can produce cleaner summaries for telehealth nurses, clinic
reception teams, and primary-care doctors. Instead of receiving a long
unstructured complaint, the clinician can see the affected region, duration,
severity, possible red flags, AI-generated follow-up questions, and the user's
own symptom notes. This can reduce repeated clarification work and make the
first human consultation more focused.

For education, the same workflow becomes a learning tool. Students can explore
anatomical regions on a 3D model, compare symptoms across body systems, and
observe how different risk factors change the recommended triage pathway. This
creates a second impact channel beyond consumer health: interactive clinical
reasoning practice for medical and nursing students, first-aid training
programmes, and public health-literacy courses.

% =============================================================================
\subsection{Market Potential and TAM--SAM--SOM Framing}
\label{sec:impact-market}
% =============================================================================

At Round~1 stage, exact commercial market sizing is less reliable than a
transparent bottom-up model. BodyCheck therefore uses a conservative
TAM--SAM--SOM framing based on reachable digital users, expected symptom-help
seeking behaviour, and the team's realistic first-year distribution channels.

\begin{table}[htbp]
  \centering
  \caption{Conservative TAM--SAM--SOM framing for BodyCheck. Values are
    planning estimates for product direction, not claimed realised traction.}
  \label{tab:tam_sam_som}
  \setlength{\tabcolsep}{5pt}
  \renewcommand{\arraystretch}{1.45}
  \small
  \begin{tabular}{@{} p{2.0cm} p{5.2cm} p{5.1cm} @{}}
    \toprule
    \textbf{Layer} & \textbf{Definition} & \textbf{BodyCheck Estimate} \\
    \midrule

    TAM &
    Digitally reachable users who may need symptom guidance, health
    information, or basic triage support. &
    Domestic floor: Vietnam's internet-user base already identified in
    \S\ref{sec:problem-fit}. Global upside comes from the English-first
    interface and web/PWA distribution. \\

    SAM &
    Users and organisations reachable by the initial product scope: English
    and Vietnamese web/mobile users, telehealth intake teams, clinics, and
    education providers. &
    Initial serviceable segment: students, young families, health-conscious
    consumers, HCMUS/HCMC beta networks, online first-aid communities, and
    small healthcare/education partners. \\

    SOM &
    Share realistically obtainable in the first 12 months after the
    competition, assuming a student-led team and low marketing budget. &
    Target: 30--50 closed-beta testers, 10{,}000--30{,}000 first-year
    consumer users, 2--3 education pilots, and 1--2 healthcare workflow
    pilots by June~2027. \\

    \bottomrule
  \end{tabular}
\end{table}

The strongest initial market is not the full healthcare system. It is the
pre-clinical decision layer: people who feel symptoms and need immediate,
understandable next-step guidance before they reach a professional. This layer
is large because it sits upstream of clinics, hospitals, telehealth platforms,
and search engines. BodyCheck's web-first architecture lets the team test this
market with a low-friction PWA before investing in heavier clinical
integrations.

The education segment provides an additional serviceable market with lower
regulatory friction. A 3D anatomical symptom checker can be packaged as a
simulation environment for instructors, allowing students to practise
triage-style reasoning without involving real patient data. This makes the
product useful even before it becomes a clinically reviewed public health tool.

% =============================================================================
\subsection{Competitive Advantage}
\label{sec:impact-competitive}
% =============================================================================

BodyCheck's competitive advantage comes from combining several capabilities
that usually appear separately in existing tools. Search engines provide
information but no structured triage path. Static symptom checkers provide
forms but often feel rigid and text-heavy. Generic chatbots can converse, but
they do not know where the user is pointing on the body and are not designed
around care-navigation handoff. Telehealth platforms connect users to
clinicians, but they often begin after the user has already decided to seek
care.

BodyCheck differs by making the body region the first-class input. The user can
click or paint affected areas on an interactive 3D anatomical model before
writing symptoms. This reduces ambiguity for non-specialist users and creates a
more structured clinical context for the AI layer. The resulting response is
not just a chat answer; it is a triage packet containing body region, symptom
summary, risk level, possible conditions, warning signs, references, and nearby
care options.

\begin{table}[htbp]
  \centering
  \caption{BodyCheck differentiation against common alternatives.}
  \label{tab:competitive_comparison}
  \setlength{\tabcolsep}{5pt}
  \renewcommand{\arraystretch}{1.45}
  \small
  \begin{tabular}{@{} p{3.2cm} p{4.8cm} p{4.8cm} @{}}
    \toprule
    \textbf{Alternative} & \textbf{Limitation} & \textbf{BodyCheck Difference} \\
    \midrule

    Search Engines &
    Return many documents; users must interpret urgency themselves. &
    Produces a constrained triage summary with next-step timing and red-flag
    escalation. \\

    Static Symptom Checkers &
    Depend on rigid forms and predefined symptom trees. &
    Supports natural language, voice input, and targeted follow-up questions. \\

    Generic AI Chatbots &
    Lack anatomical context, care-routing workflow, and product-level safety
    framing. &
    Combines 3D body-region input, structured diagnosis schema, references,
    disclaimer, and clinic handoff. \\

    Telehealth Intake Forms &
    Often start after the user has already chosen to seek care. &
    Supports the earlier decision layer and can still export structured
    summaries into telehealth workflows. \\

    \bottomrule
  \end{tabular}
\end{table}

This combination creates more than a cosmetic difference. The core workflow is
at least 30\% different from conventional text-only symptom tools because the
product changes the input modality, response structure, and downstream action
path. It begins with spatial symptom localisation, adds LLM-based reasoning,
preserves safety constraints, and ends with either self-care guidance,
professional escalation, or a clinic/education handoff.

% =============================================================================
\subsection{Revenue and Value Model}
\label{sec:impact-revenue}
% =============================================================================

BodyCheck's revenue model is designed around three complementary channels so
that the product is not dependent on a single customer type.

\begin{enumerate}[leftmargin=2em, itemsep=4pt, topsep=4pt]
  \item \textbf{B2C freemium consumer app:}
        Basic triage guidance remains free to maximise social reach. Paid
        features can include saved symptom history, voice accessibility,
        advanced source summaries, care-navigation reminders, and family
        profile management.
  \item \textbf{B2B2C healthcare API licensing:}
        Telehealth operators, clinics, hospital intake portals, and insurance
        triage teams can license BodyCheck as a white-label REST API. Pricing
        combines a platform fee with per-request usage, following the same
        linear cost logic defined in the feasibility section.
  \item \textbf{Education institutional licence:}
        Universities, medical training programmes, first-aid organisations,
        and online learning platforms can use BodyCheck as a white-label
        simulation tool. This channel is priced as an annual licence because
        value is driven by curriculum usage rather than individual medical
        sessions.
\end{enumerate}

The value model differs by channel. For consumers, the value is clarity and
confidence in the next step. For healthcare partners, the value is a more
structured pre-consultation intake flow and lower repetitive clarification
work. For education partners, the value is interactive, repeatable anatomy and
triage training without exposing real patient data.

% =============================================================================
\subsection{Impact Measurement Plan}
\label{sec:impact-measurement}
% =============================================================================

The post-competition roadmap in \S\ref{sec:feasibility-gtm} becomes stronger
when paired with measurable impact indicators. BodyCheck will track impact
through anonymised product and workflow metrics rather than raw symptom logs:

\begin{itemize}[leftmargin=2em, itemsep=3pt, topsep=4pt]
  \item \textbf{Access:} number of completed triage sessions, active users,
        returning users, and device/platform distribution.
  \item \textbf{Clarity:} post-session rating on whether the user understood
        what to do next.
  \item \textbf{Safety:} number of red-flag escalations, emergency-care
        recommendation displays, and disclaimer visibility checks.
  \item \textbf{Workflow value:} number of generated intake summaries,
        clinician/partner feedback on summary usefulness, and reduction in
        repeated clarification questions during pilot workflows.
  \item \textbf{Education value:} number of completed training cases,
        instructor feedback, and learner quiz-score improvement before and
        after using the simulation flow.
\end{itemize}

These metrics allow the team to demonstrate social benefit without making
unverifiable medical-outcome claims. The goal for the first 12 months is not
to prove that BodyCheck replaces clinical care. The goal is to prove that the
product helps users understand urgency earlier, helps care teams receive
cleaner symptom context, and helps educators teach anatomical symptom
recognition in a more interactive way.

% =============================================================================
% End of Section: Projected Impact Assessment
% =============================================================================
.
%  Uses only packages and colour definitions already loaded in main.tex.
% =============================================================================

\section{Projected Impact Assessment}
\label{sec:impact}

This section addresses the fourth Round~1 evaluation criterion ---
Projected Impact. We evaluate BodyCheck along four dimensions required by
the rubric: social and business value; potential user scale using a
TAM--SAM--SOM framing; competitive advantage in the target market; and the
revenue/value model that can sustain the product after the competition.

The impact thesis is straightforward: BodyCheck does not try to replace
doctors. Its highest-value contribution is to improve the earliest minutes of
a health concern, when users are uncertain, anxious, and often forced to choose
between ignoring symptoms, searching online, or visiting a clinic without a
clear sense of urgency. By turning vague symptoms into structured triage
guidance, BodyCheck can reduce confusion for users, prepare cleaner intake
summaries for healthcare teams, and create a reusable learning environment for
health-literacy and medical-training use cases.

Table~\ref{tab:impact_matrix} summarises the expected impact channels.

% ---------------------------------------------------------------------------
% Table 1 -- Projected Impact Matrix
% ---------------------------------------------------------------------------
\begin{table}[htbp]
  \centering
  \caption{Projected impact channels for BodyCheck aligned to the
    HackAIthon 2026 Round~1 impact rubric.}
  \label{tab:impact_matrix}
  \setlength{\tabcolsep}{5pt}
  \renewcommand{\arraystretch}{1.45}
  \small
  \begin{tabular}{@{} p{3.0cm} p{4.6cm} p{4.7cm} @{}}
    \toprule
    \textbf{Impact Area} &
    \textbf{Expected Benefit} &
    \textbf{Measurement Signal} \\
    \midrule

    Social Benefit &
    Earlier recognition of red flags, clearer self-care guidance, and
    faster escalation to professional care when symptoms appear risky. &
    Time-to-first-guidance, emergency-escalation rate, user clarity score,
    and clinic handoff clicks. \\

    Healthcare Operations &
    Cleaner pre-consultation summaries reduce repetitive intake work for
    triage operators, telehealth nurses, and primary-care clinics. &
    Intake completion rate, summary reuse rate, average number of follow-up
    clarification questions. \\

    Education &
    Interactive 3D body-region selection and AI-generated triage cases create
    a simulation tool for medical, nursing, first-aid, and health-literacy
    programmes. &
    Number of institutional pilots, completed learning cases, instructor
    feedback score. \\

    Business Value &
    One shared technical core supports three commercial channels: consumer
    app, healthcare API, and education licence. &
    Monthly active users, API request volume, institutional pilots, paid
    conversion rate. \\

    \bottomrule
  \end{tabular}
\end{table}

% =============================================================================
\subsection{Social and Healthcare Impact}
\label{sec:impact-social}
% =============================================================================

BodyCheck creates value by reducing the gap between symptom onset and an
appropriate next action. For non-specialist users, the product translates a
stressful and ambiguous experience --- for example, ``chest tightness'',
``shoulder pain'', or ``lower abdominal discomfort'' --- into a structured
triage response:

\begin{itemize}[leftmargin=2em, itemsep=3pt, topsep=4pt]
  \item possible explanations expressed with uncertainty language;
  \item risk level and recommended timing for care;
  \item red-flag symptoms that require urgent escalation;
  \item home-care tips for low-risk situations;
  \item trusted references and nearby care options when professional help is
        appropriate.
\end{itemize}

This matters because the first decision is often the most confusing one:
whether to wait, self-care, schedule a routine appointment, or seek urgent
care. BodyCheck improves that decision without making a final diagnosis. The
AI response is deliberately constrained by prompt rules, a structured schema,
and a medical disclaimer so that the user receives triage support rather than
false diagnostic certainty.

For healthcare operations, the impact is not only patient-facing. A structured
intake flow can produce cleaner summaries for telehealth nurses, clinic
reception teams, and primary-care doctors. Instead of receiving a long
unstructured complaint, the clinician can see the affected region, duration,
severity, possible red flags, AI-generated follow-up questions, and the user's
own symptom notes. This can reduce repeated clarification work and make the
first human consultation more focused.

For education, the same workflow becomes a learning tool. Students can explore
anatomical regions on a 3D model, compare symptoms across body systems, and
observe how different risk factors change the recommended triage pathway. This
creates a second impact channel beyond consumer health: interactive clinical
reasoning practice for medical and nursing students, first-aid training
programmes, and public health-literacy courses.

% =============================================================================
\subsection{Market Potential and TAM--SAM--SOM Framing}
\label{sec:impact-market}
% =============================================================================

At Round~1 stage, exact commercial market sizing is less reliable than a
transparent bottom-up model. BodyCheck therefore uses a conservative
TAM--SAM--SOM framing based on reachable digital users, expected symptom-help
seeking behaviour, and the team's realistic first-year distribution channels.

\begin{table}[htbp]
  \centering
  \caption{Conservative TAM--SAM--SOM framing for BodyCheck. Values are
    planning estimates for product direction, not claimed realised traction.}
  \label{tab:tam_sam_som}
  \setlength{\tabcolsep}{5pt}
  \renewcommand{\arraystretch}{1.45}
  \small
  \begin{tabular}{@{} p{2.0cm} p{5.2cm} p{5.1cm} @{}}
    \toprule
    \textbf{Layer} & \textbf{Definition} & \textbf{BodyCheck Estimate} \\
    \midrule

    TAM &
    Digitally reachable users who may need symptom guidance, health
    information, or basic triage support. &
    Domestic floor: Vietnam's internet-user base already identified in
    \S\ref{sec:problem-fit}. Global upside comes from the English-first
    interface and web/PWA distribution. \\

    SAM &
    Users and organisations reachable by the initial product scope: English
    and Vietnamese web/mobile users, telehealth intake teams, clinics, and
    education providers. &
    Initial serviceable segment: students, young families, health-conscious
    consumers, HCMUS/HCMC beta networks, online first-aid communities, and
    small healthcare/education partners. \\

    SOM &
    Share realistically obtainable in the first 12 months after the
    competition, assuming a student-led team and low marketing budget. &
    Target: 30--50 closed-beta testers, 10{,}000--30{,}000 first-year
    consumer users, 2--3 education pilots, and 1--2 healthcare workflow
    pilots by June~2027. \\

    \bottomrule
  \end{tabular}
\end{table}

The strongest initial market is not the full healthcare system. It is the
pre-clinical decision layer: people who feel symptoms and need immediate,
understandable next-step guidance before they reach a professional. This layer
is large because it sits upstream of clinics, hospitals, telehealth platforms,
and search engines. BodyCheck's web-first architecture lets the team test this
market with a low-friction PWA before investing in heavier clinical
integrations.

The education segment provides an additional serviceable market with lower
regulatory friction. A 3D anatomical symptom checker can be packaged as a
simulation environment for instructors, allowing students to practise
triage-style reasoning without involving real patient data. This makes the
product useful even before it becomes a clinically reviewed public health tool.

% =============================================================================
\subsection{Competitive Advantage}
\label{sec:impact-competitive}
% =============================================================================

BodyCheck's competitive advantage comes from combining several capabilities
that usually appear separately in existing tools. Search engines provide
information but no structured triage path. Static symptom checkers provide
forms but often feel rigid and text-heavy. Generic chatbots can converse, but
they do not know where the user is pointing on the body and are not designed
around care-navigation handoff. Telehealth platforms connect users to
clinicians, but they often begin after the user has already decided to seek
care.

BodyCheck differs by making the body region the first-class input. The user can
click or paint affected areas on an interactive 3D anatomical model before
writing symptoms. This reduces ambiguity for non-specialist users and creates a
more structured clinical context for the AI layer. The resulting response is
not just a chat answer; it is a triage packet containing body region, symptom
summary, risk level, possible conditions, warning signs, references, and nearby
care options.

\begin{table}[htbp]
  \centering
  \caption{BodyCheck differentiation against common alternatives.}
  \label{tab:competitive_comparison}
  \setlength{\tabcolsep}{5pt}
  \renewcommand{\arraystretch}{1.45}
  \small
  \begin{tabular}{@{} p{3.2cm} p{4.8cm} p{4.8cm} @{}}
    \toprule
    \textbf{Alternative} & \textbf{Limitation} & \textbf{BodyCheck Difference} \\
    \midrule

    Search Engines &
    Return many documents; users must interpret urgency themselves. &
    Produces a constrained triage summary with next-step timing and red-flag
    escalation. \\

    Static Symptom Checkers &
    Depend on rigid forms and predefined symptom trees. &
    Supports natural language, voice input, and targeted follow-up questions. \\

    Generic AI Chatbots &
    Lack anatomical context, care-routing workflow, and product-level safety
    framing. &
    Combines 3D body-region input, structured diagnosis schema, references,
    disclaimer, and clinic handoff. \\

    Telehealth Intake Forms &
    Often start after the user has already chosen to seek care. &
    Supports the earlier decision layer and can still export structured
    summaries into telehealth workflows. \\

    \bottomrule
  \end{tabular}
\end{table}

This combination creates more than a cosmetic difference. The core workflow is
at least 30\% different from conventional text-only symptom tools because the
product changes the input modality, response structure, and downstream action
path. It begins with spatial symptom localisation, adds LLM-based reasoning,
preserves safety constraints, and ends with either self-care guidance,
professional escalation, or a clinic/education handoff.

% =============================================================================
\subsection{Revenue and Value Model}
\label{sec:impact-revenue}
% =============================================================================

BodyCheck's revenue model is designed around three complementary channels so
that the product is not dependent on a single customer type.

\begin{enumerate}[leftmargin=2em, itemsep=4pt, topsep=4pt]
  \item \textbf{B2C freemium consumer app:}
        Basic triage guidance remains free to maximise social reach. Paid
        features can include saved symptom history, voice accessibility,
        advanced source summaries, care-navigation reminders, and family
        profile management.
  \item \textbf{B2B2C healthcare API licensing:}
        Telehealth operators, clinics, hospital intake portals, and insurance
        triage teams can license BodyCheck as a white-label REST API. Pricing
        combines a platform fee with per-request usage, following the same
        linear cost logic defined in the feasibility section.
  \item \textbf{Education institutional licence:}
        Universities, medical training programmes, first-aid organisations,
        and online learning platforms can use BodyCheck as a white-label
        simulation tool. This channel is priced as an annual licence because
        value is driven by curriculum usage rather than individual medical
        sessions.
\end{enumerate}

The value model differs by channel. For consumers, the value is clarity and
confidence in the next step. For healthcare partners, the value is a more
structured pre-consultation intake flow and lower repetitive clarification
work. For education partners, the value is interactive, repeatable anatomy and
triage training without exposing real patient data.

% =============================================================================
\subsection{Impact Measurement Plan}
\label{sec:impact-measurement}
% =============================================================================

The post-competition roadmap in \S\ref{sec:feasibility-gtm} becomes stronger
when paired with measurable impact indicators. BodyCheck will track impact
through anonymised product and workflow metrics rather than raw symptom logs:

\begin{itemize}[leftmargin=2em, itemsep=3pt, topsep=4pt]
  \item \textbf{Access:} number of completed triage sessions, active users,
        returning users, and device/platform distribution.
  \item \textbf{Clarity:} post-session rating on whether the user understood
        what to do next.
  \item \textbf{Safety:} number of red-flag escalations, emergency-care
        recommendation displays, and disclaimer visibility checks.
  \item \textbf{Workflow value:} number of generated intake summaries,
        clinician/partner feedback on summary usefulness, and reduction in
        repeated clarification questions during pilot workflows.
  \item \textbf{Education value:} number of completed training cases,
        instructor feedback, and learner quiz-score improvement before and
        after using the simulation flow.
\end{itemize}

These metrics allow the team to demonstrate social benefit without making
unverifiable medical-outcome claims. The goal for the first 12 months is not
to prove that BodyCheck replaces clinical care. The goal is to prove that the
product helps users understand urgency earlier, helps care teams receive
cleaner symptom context, and helps educators teach anatomical symptom
recognition in a more interactive way.

% =============================================================================
% End of Section: Projected Impact Assessment
% =============================================================================
